\section{Comparative results}\label{sec:results}

\subsection{Development history and fresh validation}
The initially locked symmetric formulation failed held-out SUMO C3 testing. A regularised revision had favourable estimates in a 30-episode replication, but both paired intervals included zero. A separately registered 160-episode confirmation again did not establish superiority; comparator advantages were approximately $0.012\,\ms$ after unit conversion. These unsuccessful stages are retained rather than reclassified as successful evidence for the final method.

The subsequent signed-tail candidate was selected using split-safe tuning and locked before its fresh validation. Table~\ref{tab:primary} reports the six final FAR-GW dataset/scenario cells. The 40-episode SUMO C3 validation passes the complete rule. SUMO C4 has positive paired score effects but coverage 0.8793, narrowly below the registered floor. Neither favourable score estimates alone nor rounding coverage to two decimals changes that decision.

METR-LA passes in both C3 and C4. PEMS-BAY C3 has small effects whose intervals include zero; PEMS-BAY C4 favours both comparators. These negative transfer results are principal findings. The method is left locked rather than retuned using the evaluated test outcomes.

\subsection{Cross-backbone and relational comparisons}
Figure~\ref{fig:benchmarks} presents every evaluated real-data comparator with the proposed method as the zero reference. Within each backbone, all calibrators use the same forecasts and masks. The plot therefore compares calibration conditional on a backbone; it is not a ranking of the predictors themselves. CoRel points appear only for FAR-GW because that is the evaluated scope. Table~\ref{tab:corel} reports the direct CoRel comparison in numerical form, including both methods' coverage and the paired score interval. All four CoRel-minus-FAR-Cal intervals are positive, but this does not overturn PEMS-BAY failures against rolling and ACI. Different comparators support different conclusions.

The four-backbone matrix passes in 12 of 16 registered cells: all eight METR-LA cells, and the four PEMS-BAY cells using STAEformer or DCRNN. FAR-GW and STID each fail both PEMS-BAY conditions. Table~\ref{tab:backbones} provides both paired score effects, coverage and joint decisions for every cell. Its effect columns compare calibrators within a backbone; they must not be read as point-model rankings. Consistency across METR-LA predictors supports transfer on that network, while PEMS-BAY demonstrates predictor-dependent behaviour.

CoRel-minus-\FAR{} differences are positive in all four FAR-GW real-data cells. After SI conversion, the estimates are 0.758, 0.973, 1.008 and 1.004\,$\ms$ for METR-LA C3/C4 and PEMS-BAY C3/C4. Their respective intervals are [0.518, 1.030], [0.616, 1.403], [0.814, 1.215] and [0.787, 1.227]\,$\ms$. CoRel coverage is 0.8743, 0.8725, 0.8837 and 0.8743. This direct comparison supports an advantage over this adaptation; it does not imply that \FAR{} beats rolling or ACI in the PEMS-BAY cells.

\subsection{Coverage and score must be read jointly}
Figure~\ref{fig:quality} places each final FAR-GW cell's coverage beside its interval score, using identical axes within the quality panel. It makes the SUMO C4 decision visible: positive score improvements coexist with coverage below 0.88. PEMS-BAY coverage exceeds the floor, yet superiority against rolling and ACI is not established. These are different failure modes. Table~\ref{tab:primary} therefore reports the conjunction, while the figure preserves each metric separately. The vertical separation of network scores is descriptive and not a between-network treatment effect.


\subsection{Core faults, independent controls and one-factor stress}
The core and matched-control matrices complete 168 and 120 stages. Table~\ref{tab:controls} reports their paired effects and joint decisions. METR-LA passes C1, C2 and C5 and both independent C3/C4 controls; PEMS-BAY passes none. The independent-control results preserve the network-specific pattern without proving that spatial correlation itself causes the method's advantage.

All 1,488 one-factor stress stages finish, yielding 48 dataset/scenario/variant summaries. METR-LA passes all 24 tested cells and PEMS-BAY none of its 24; every aggregate coverage is at least 0.88. Table~\ref{tab:stress} groups the findings by factor, and Figure~\ref{fig:duration} shows score sensitivity to outage duration. Duration, connected fraction, backlog delay, permanent loss and initial archive impairment test different limits of the reporting system; they are not pooled into a single generalisation guarantee. The MNAR grid retains the same distinction: favourable performance in one network and failure in the other remain compatible with selective outcome observation.


\subsection{Additional reporting stress}
Figure~\ref{fig:stress} displays affected fraction, extra backlog delay, permanent-loss probability, cold calibration and low-speed MNAR sensitivity. Each point compares ACI with FAR-Cal using the registered paired inference, with positive values favouring FAR-Cal. Continuous-factor lines connect tested values only; cold calibration and MNAR are separate categorical experiments. All METR-LA stress cells pass the joint rule, whereas no PEMS-BAY stress cell does. Table~\ref{tab:stress} groups these decisions; the plotted ACI contrast alone is not the full two-comparator rule. The cold archive tests whether an initially impaired calibration source changes performance, and the MNAR panel tests selective outcome suppression without claiming correction.

The duration curve in Figure~\ref{fig:duration} reports absolute FAR-Cal score rather than a paired treatment effect. Longer outages can change which targets enter the fault-focused audit, so the curve cannot isolate a duration effect on a fixed target population. It complements the paired stress results by showing the magnitude of interval-score penalties in each tested condition. These plots are based on completed evaluations, not simulated illustrative values or interpolated predictions of untested fault levels.
